% from ega2-1-fr.tex
\oldpage[II]{191}
(\hyperref[II.8.8.5.2-fr]{8.8.5.2}) et
(\hyperref[II.8.9.1.1-fr]{8.9.1.1}) montre que la restriction de $u$ au
sous-préschéma fermé $\varepsilon'(Y)$ de $C'$ est un isomorphisme sur le
sous-préschéma fermé $\varepsilon(Y)$ de $C$, d'où résulte aussitôt que pour
tout $z'\in\varepsilon'(Y)$, si $z=u(z')$, $u$ définit un isomorphisme de
$k(z)$ sur $k(z')$. Ces remarques prouvent que $u$ est bijectif et radiciel~;
en outre, si $\psi$ et $\psi'$ sont les morphismes structuraux $C\to Y$,
$C'\to Y$, on a $\psi'=\psi\circ u$ et comme $\psi'$ est séparé
(\hyperref[II.1.2.4-fr]{1.2.4}), il en est de même de $u$
(\hyperref[I.5.5.1-fr]{I, 5.5.1, (v)}). On a vu d'autre part dans la
démonstration du lemme (\hyperref[II.8.11.5-fr]{8.11.5}) que $q'$ est
surjectif~; comme $q=u\circ q'$ est propre, on conclut finalement de
(\hyperref[II.5.4.3-fr]{5.4.3}) et (\hyperref[II.5.4.9-fr]{5.4.9}) que $u$
est universellement fermé.
\end{proof}

\begin{proposition}[8.11.7]
\label{II.8.11.7-fr}
Soient $Y$ un préschéma, $X$ un préschéma intègre, $p:X\to Y$ un morphisme
propre, $\mathcal{L}$ un $\mathcal{O}_X$-Module inversible $p$-ample, $C$ un
$Y$-préschéma normal, $\varepsilon:Y\to C$ une $Y$-section,
$q:V=\mathbf{V}(\mathcal{L})\to C$ un $Y$-morphisme, tels que le diagramme
(\hyperref[II.8.9.1.1-fr]{8.9.1.1}) soit commutatif. On suppose en outre que,
si $p=(p_0,\theta)$, $\theta:\mathcal{O}_Y\to p_*(\mathcal{O}_X)$ soit un
isomorphisme. Soient alors
$\mathcal{S}'=\bigoplus_{n\geqslant0}p_*(\mathcal{L}^{\otimes n})$,
$C'=\operatorname{Spec}(\mathcal{S}')$, et
$q':\mathbf{V}(\mathcal{L})\to C'$ le $Y$-morphisme canonique
(\hyperref[II.8.8.5-fr]{8.8.5}). Alors l'unique $Y$-morphisme $u:C'\to C$ tel
que $q=u\circ q'$ est un isomorphisme.
\end{proposition}

\begin{proof}
Il résulte de (\hyperref[II.8.8.6-fr]{8.8.6}) que $C'$ est intègre~; comme
$u$ est un homéomorphisme des espaces sous-jacents $C'\to C$ ($u$ étant
bijectif et fermé d'après (\hyperref[II.8.11.6-fr]{8.11.6})), $C$ est
irréductible, donc intègre, et puisque la restriction de $u$ à un ouvert non
vide de $C'$ est un isomorphisme sur un ouvert de $C$, $u$ est birationnel.
Puisque $C$ est supposé normal, il suffit d'appliquer
(\hyperref[II.8.11.2-fr]{8.11.2}) pour obtenir la conclusion.
\end{proof}

\begin{remark}[8.11.8]
\label{II.8.11.8-fr}
\begin{enumerate}
  \item[\rm{(i)}] Notons que l'hypothèse que $C$ est normal implique qu'il en
  est de même de $X$. En effet, $C'=\operatorname{Spec}(\mathcal{S}')$ est
  alors normal, étant isomorphe à $C$, et intègre en vertu de
  (\hyperref[II.8.8.6-fr]{8.8.6})~; on en conclut que
  $\operatorname{Proj}(\mathcal{S}')$ est \emph{normal}. En effet, la question
  est locale sur $Y$~; si $Y$ est affine, $\mathcal{S}'=\widetilde{S'}$,
  l'anneau $S'=\Gamma(C',\mathcal{O}_{C'})$ est intègre et intégralement clos
  (\hyperref[II.8.8.6.1-fr]{8.8.6.1}), donc, pour tout élément homogène
  $f\in S'_+$, l'anneau gradué $S'_f$ est intègre et intégralement clos
  ([13], t.~I, p.~257 et 261), et par suite aussi l'anneau $S'_{(f)}$ de ses
  termes de degré 0, car l'intersection de $S'_f$ et du corps des fractions de
  $S'_{(f)}$ est égale à $S'_{(f)}$~; ce qui prouve notre assertion
  (\hyperref[II.6.3.4-fr]{6.3.4}). Enfin, comme $X$ est isomorphe à un
  sous-préschéma ouvert de $\operatorname{Proj}(\mathcal{S}')$
  (\hyperref[II.8.8.1-fr]{8.8.1}), $X$ est bien normal. On peut donc exprimer
  la prop. (\hyperref[II.8.11.7-fr]{8.11.7}) sous la forme suivante~:
  \emph{Si $X$ est intègre et normal,
  $p=(p_0,\theta):X\to Y$ un morphisme propre tel que
  $\theta:\mathcal{O}_Y\to p_*(\mathcal{O}_X)$ soit un isomorphisme, alors,
  pour tout $\mathcal{O}_X$-Module $p$-ample $\mathcal{L}$, il existe une
  façon et une seule de contracter la section nulle de
  $V=\mathbf{V}(\mathcal{L})$ de façon à obtenir un $Y$-schéma normal $C$ et
  un $Y$-morphisme propre $q:V\to C$.}
  \item[\rm{(ii)}] Lorsque $p$ est propre, l'hypothèse
  $p_*(\mathcal{O}_X)=\mathcal{O}_Y$ peut être considérée comme une hypothèse
  auxiliaire, ne constituant pas une vraie restriction de la généralité du
  résultat. En effet, si elle n'est pas vérifiée, il suffit de remplacer $Y$
  par le $Y$-schéma
  $Y'=\operatorname{Spec}(p_*(\mathcal{O}_X))$, et de considérer $X$ comme un
  $Y'$-schéma. Nous reviendrons sur cette méthode générale au chapitre III,
  \S~4.
\end{enumerate}
\end{remark}
\subsection{Faisceaux quasi-cohérents sur les cônes projetants}
\label{subsection:II.8.12-fr}

\begin{env}[8.12.1]
\label{II.8.12.1-fr}
Reprenons les hypothèses et notations de
(\hyperref[II.8.3.1-fr]{8.3.1}). Soit $\mathcal{M}$ un
$\mathcal{S}$-Module gradué quasi-cohérent~; pour éviter toute confusion,
nous désignerons par $\widetilde{\mathcal{M}}$ le $\mathcal{O}_C$-Module

